% Generated draft: requires mathematical and source review.
\documentclass[11pt]{article}
\usepackage[margin=25mm]{geometry}
\usepackage{fontspec}
\setmainfont{Latin Modern Roman}
\usepackage{amsmath,amssymb,mathtools,unicode-math}
\setmathfont{Latin Modern Math}
\usepackage{enumitem,needspace,xurl,hyperref}
\setlength{\parindent}{0pt}
\setlength{\parskip}{5pt}
\setlength{\emergencystretch}{3em}
\newcommand{\N}{\mathbb{N}}
\newcommand{\Z}{\mathbb{Z}}
\newcommand{\Q}{\mathbb{Q}}
\newcommand{\R}{\mathbb{R}}
\newcommand{\C}{\mathbb{C}}
\newcommand{\F}{\mathbb{F}}
\newcommand{\A}{\mathbb{A}}
\newcommand{\PP}{\mathbb{P}}
\newcommand{\E}{\mathbb{E}}
\newcommand{\eps}{\varepsilon}
\newcommand{\dd}{\,\mathrm{d}}
\begin{document}
\section*{UnsolvedMath ID: 3012}
\textbf{Source problem code:} catalogue entry
\subsection{Definitions and mathematical statement}
Consider topological manifold pairs $(M_i^{n+k},N_i^n)$, $i=1,2$, with locally flat embeddings $N_i\subset M_i$. At an interior point, local flatness means that a neighborhood pair is homeomorphic to $(\R^{n+k},\R^n)$ with the standard linear inclusion. In the usual positive-dimensional connected-sum construction, choose embedded standard ball pairs $(D^{n+k},D^n)$ around points of $N_i$, remove their interiors, and glue the boundary pairs
\[
 (S^{n+k-1},S^{n-1})
\]
by a compatible homeomorphism. This produces a candidate pair denoted $(M_1,N_1)\#(M_2,N_2)$. If orientations are imposed, the charts and gluing must respect the selected connected-sum orientation convention.

\textbf{Imported choice-independence question.} Do different compatible choices of points, local ball-pair charts and gluing data always produce equivalent topological pairs? In a homeomorphism-of-pairs version, equivalence means a homeomorphism between the ambient manifolds carrying one embedded submanifold onto the other, preserving orientations when those are part of the data. The source does not say whether it instead asks for ambient isotopy, a relative/marked equivalence, or un-oriented homeomorphism. It also omits connectedness, orientations, local-orientation conventions, dimension range and possible boundary behavior. These choices are necessary to specify exactly which universal choice-independence assertion is intended.

The literature triage dated 2026-08-17 reports an affirmative codimension-one case from Brown's local-flatness results. It reports Livingston's positive result in codimension two in every dimension \emph{for connected, oriented, locally flat pairs}, with independence up to orientation-preserving homeomorphism. Those hypotheses and that equivalence notion are retained, rather than extending the result to all pairs left unspecified in the short statement.

The source explains that normal-bundle uniqueness is essential in codimension two and fails in higher codimensions. It leaves codimension $k>2$ with $n+k>4$ unresolved: prove the appropriately specified connected sum independent of its data, or find pairs where it is not. It suggests that a pairwise torus trick could yield a solution if a corresponding pairwise version of Wall's non-simply-connected surgery were known and in print; this is a conditional proposed route, not a theorem supplied here.

The registry label is open and the triage label partially solved, reflecting the restricted low-codimension results. No source, PDF or topological proof was independently checked, and no present-day status assessment is claimed.
\subsection{Short English statement}
Remove standard ball pairs from two locally flat manifold pairs and glue their boundary pairs. Does the resulting topological pair depend on the choices? The supplied review reports codimension-one results and a connected oriented codimension-two theorem, while leaving higher codimension in ambient dimension above four open. The exact equivalence and orientation conventions still need specification.
\subsection{Sources}
\begin{itemize}
\item UnsolvedMath local registry, data/unsolvedmath/problems.json, permanent ID 3012, source code catalogue entry, ``Kirby Problem 5.5,'' with the pair connected-sum question and triage dated 2026-08-17.
\item R. I. Baykur, R. C. Kirby and D. Ruberman, \emph{K3: A New Problem List in Low-Dimensional Topology}, Problem 5.5, author's preliminary AMS version (2026), \url{https://bpb-us-e2.wpmucdn.com/websites.umass.edu/dist/b/22144/files/2026/04/K3-problem-list-watermarked.pdf}. The PDF was not opened.
\item [Bro62] Morton Brown, ``Locally flat imbeddings of topological manifolds,'' \emph{Ann. of Math. (2)} 75 (1962), 331--341, DOI 10.2307/1970177, supplied review link \url{https://doi.org/10.2307/1970177}.
\item [Liv24] Charles Livingston, ``Connected sums of codimension two locally flat submanifolds,'' \emph{Proc. Roy. Soc. Edinburgh Sect. A} 154(6) (2024), 1937--1944, DOI 10.1017/prm.2022.87, supplied review link \url{https://doi.org/10.1017/prm.2022.87}; the background points to Appendix C for normal-bundle issues.
\item [KS77] Robion C. Kirby and Laurence C. Siebenmann, \emph{Foundational essays on topological manifolds, smoothings, and triangulations}, Annals of Mathematics Studies 88, Princeton University Press (1977), with notes by John Milnor and Michael Atiyah.
\item [Wal99] C. T. C. Wall, \emph{Surgery on compact manifolds}, Mathematical Surveys and Monographs 69, AMS, second edition (1999), edited and with foreword by A. A. Ranicki, DOI 10.1090/surv/069.
\item Ulam AI and the UnsolvedMath Contributors, \emph{UnsolvedMath: A Curated Collection of Open Mathematics Problems}. Dataset curation and original metadata are supplied under CC BY 4.0; referenced source material retains its own terms. No external reference was independently consulted for this draft.
\end{itemize}
\end{document}
